%! Operations with Polynomials — Unit 3, Lesson 3.1 — Exit Ticket (Answer Key)
\documentclass[10pt]{article}
\usepackage{algebra2-article}
\usepackage{algebra2-key}   % pulls in -boxes; do NOT also load -boxes
\newcommand{\TallMath}[1]{$\displaystyle #1\rule[-1.4em]{0pt}{3.2em}$}

\begin{document}

\pageheader{Unit 3, Lesson 3.1}{Exit Ticket --- Answer Key}
\namedateperiod

\vspace{0.10in}

No notes. Write every answer in \textbf{standard form} and show your work.

\begin{enumerate}[leftmargin=1.8em, itemsep=12pt]
  \item \textbf{Subtract.} $(4x^3 - 2x^2 + 7) - (x^3 + 5x^2 - 3x + 7)$

    \vspace{3pt}
    \emph{First rewrite with every sign changed:} \ans{$4x^3 - 2x^2 + 7 - x^3 - 5x^2 + 3x - 7$}

    \vspace{5pt}
    \emph{Answer:} \ans{$3x^3 - 7x^2 + 3x$}

  \item \textbf{Multiply, then describe.} $(2x - 3)(x^2 + 4x - 5)$
    \begin{center}
    \renewcommand{\arraystretch}{1.9}
    \begin{tabular}{|c|c|c|c|}
    \hline
    $\times$ & $x^2$ & $+4x$ & $-5$ \\ \hline
    $2x$ & \ans{$2x^3$} & \ans{$8x^2$} & \ans{$-10x$} \\ \hline
    $-3$ & \ans{$-3x^2$} & \ans{$-12x$} & \ans{$15$} \\ \hline
    \end{tabular}
    \end{center}
    \emph{Answer:} \ans{$2x^3 + 5x^2 - 22x + 15$}

    \vspace{3pt}
    Its degree is \ans{3} and its leading coefficient is \ans{2}, so as
    $x\to-\infty$ the graph goes \ans{down} and as $x\to+\infty$ it goes \ans{up}.

  \item \textbf{SOL-style multiple choice.} Which expression is equivalent to $(3x - 4)^2$?
    \begin{enumerate}[label=\Alph*., leftmargin=1.9em, itemsep=1pt]
      \item $9x^2 - 16$
      \item $9x^2 + 16$
      \item $9x^2 - 12x + 16$
      \item $9x^2 - 24x + 16$ \quad \textcolor{keyred}{\textbf{$\leftarrow$ correct}}
    \end{enumerate}
    \begin{itemize}[leftmargin=1.5em, nosep]
      \item $(3x-4)^2 = (3x)^2 - 2(3x)(4) + 4^2 = 9x^2 - 24x + 16$.
      \item \textbf{A} is the \emph{difference of squares} $(3x+4)(3x-4)$ --- a different problem.
      \item \textbf{B} drops the middle term entirely: the classic $(a-b)^2 = a^2+b^2$ error.
      \item \textbf{C} computes the middle term as $ab$ instead of $2ab$ --- half the right answer.
    \end{itemize}
\end{enumerate}

\begin{teachernote}
Item~1 is deliberately built so the constants cancel: an answer ending in $+14$ means only the first
sign was distributed, and an answer with no $+3x$ means the $-3x$ term was skipped because it had
``nothing to pair with.'' Item~2's last line is the Lesson~3.0 link --- students who expand correctly
but stall on the arms need end-behavior review, not multiplication review, so sort those two piles
separately. Item~3's distractors map one-to-one onto the three errors above; a class-wide spike on
\textbf{B} or \textbf{C} is worth a five-minute re-open of Lesson~3.2.
\end{teachernote}

\end{document}
